% !TEX root = ../../tir_monograph_v12.tex
\chapter{First Distinction and the Relational Kernel}
\label{ch:v12_first_distinction}

\section{Dependency before temporal order}

The monograph begins with structural dependency. Write
\begin{equation}
\boxed{X\prec_{\rm dep}Y}
\end{equation}
when the admitted structure represented by $Y$ uses $X$ as an upstream parent in the derivation graph. Temporal ordering receives its own dynamical treatment in the companion time branch; the primitive relation used here is therefore a dependency order.

The zero-distinction boundary is denoted
\begin{equation}
D_0=0.
\end{equation}
The first admitted non-empty relational carrier is the point-like minimal carrier
\begin{equation}
\boxed{\mathcal P=\text{minimal non-empty distinguishable carrier}.}
\end{equation}
This is the A1 root of the TIR axiomatic kernel.
\vTwelveStatus{B}{--}{PASS}
The status is the declared foundational role of the carrier; downstream mathematical claims are validated separately.

\section{Primitive kernel}

The eight TIR axioms are grouped by role rather than by historical order:
\begin{description}
\item[A1 -- point minimality:] supplies the minimal non-empty carrier $\mathcal P$.
\item[A2 -- quantum point:] supplies a normalized Hilbert-space representation for the minimal carrier.
\item[A3 -- information primacy:] supplies distinguishable informational relations and their normalized information carrier.
\item[A4 -- spherical geometric efficiency:] selects the sphere for the declared isotropic boundary/volume functional.
\item[A5 -- arithmetic measures geometry:] reads geometric closure through discrete invariants such as winding or multiplicity.
\item[A6 -- natural numbers from complex phase closure:] assigns natural closure indices to admitted complex phase structure.
\item[A7 -- universal symmetry:] supplies symmetry-governed relational laws and the primitive exchange action.
\item[A8 -- paradox stabilization:] lifts context-dependent incompatible projections into a higher-order closure carrier.
\end{description}

This partition creates a circularity firewall: the first binary distinction is derived from a nontrivial involution and exchange symmetry, while arithmetic closure indices enter later through A5--A6.

\section{Primitive exchange orbit}

Let $x$ be a primitive relational state and let $J$ be a nontrivial involution,
\begin{equation}
\boxed{J^2=\mathrm{id},\qquad Jx\neq x.}
\label{eq:v12_primitive_involution}
\end{equation}
The orbit generated by $J$ is
\begin{equation}
\mathcal O_J(x)=\{x,Jx\}.
\end{equation}
Since
\[
J(Jx)=x,
\]
iteration closes on the exchanged partner. Rename the two relational poles
\begin{equation}
\boxed{\mathcal O_J(x)=\{N,S\}.}
\label{eq:v12_primitive_binary_orbit}
\end{equation}

\begin{theorem}[Primitive binary orbit]
Given a state $x$ and a nontrivial involution $J$, the orbit of $x$ under the group generated by $J$ contains exactly the two distinct elements $x$ and $Jx$.
\end{theorem}

The theorem uses ordinary mathematical cardinality to describe the orbit. The TIR arithmetic-emergence branch remains downstream of this relational result.
\vTwelveStatus{B}{--}{PASS}

\section{Unique exchange-invariant share}

Attach normalized nonnegative relational weights
\begin{equation}
w_N,w_S\ge0,
\qquad
w_N+w_S=1.
\end{equation}
Exchange invariance gives
\begin{equation}
(w_N,w_S)=(w_S,w_N),
\end{equation}
therefore
\begin{equation}
w_N=w_S.
\end{equation}
Normalization fixes the unique symmetric assignment,
\begin{equation}
\boxed{
(w_N,w_S)=\left(\frac12,\frac12\right).
}
\label{eq:v12_unique_half_share}
\end{equation}

\begin{theorem}[Unique exchange-invariant share]
The unique normalized nonnegative weight assignment invariant under exchange of the two primitive poles is $(1/2,1/2)$.
\end{theorem}

Introducing the scalar coordinate
\[
u=w_S,
\qquad
w_N=1-u,
\]
turns the pole exchange into
\begin{equation}
J(u)=1-u.
\end{equation}
Its unique fixed point is
\begin{equation}
\boxed{\operatorname{Fix}(J)=\left\{\frac12\right\}.}
\label{eq:v12_half_fixed_point}
\end{equation}
Thus the normalized exchange-invariant vector and the half coordinate are two descriptions of the same primitive balance.

\section{First information value}

For the normalized binary relation, use the Shannon carrier
\begin{equation}
H_2(u)=-(1-u)\ln(1-u)-u\ln u.
\end{equation}
At the exchange-fixed share,
\begin{equation}
\boxed{H_2(1/2)=\ln2.}
\label{eq:v12_first_ln2}
\end{equation}
This yields the first exact scalar dependency chain
\begin{equation}
\boxed{
\mathcal P
\xrightarrow{J^2=1,\,Jx\neq x}
\{N,S\}
\xrightarrow{\rm exchange\ invariance}
\frac12
\xrightarrow{H_2}
\ln2.
}
\label{eq:v12_primitive_scalar_chain}
\end{equation}
Its minimal TIR parent set is A1, A3 and A7. A2 enters at the coherent quantum lift.

\section{Quantum lift}

Represent the two poles by an orthonormal basis
\[
|N\rangle,\qquad |S\rangle.
\]
Their minimal complex span is
\begin{equation}
\boxed{\mathcal H_{NS}\cong\mathbb C^2.}
\label{eq:v12_binary_c2}
\end{equation}
The balanced coherent family is
\begin{equation}
\boxed{
|\psi_{1/2}(\varphi)\rangle
=\frac{|N\rangle+e^{i\varphi}|S\rangle}{\sqrt2},
\qquad
\varphi\in\mathbb R/2\pi\mathbb Z.
}
\label{eq:v12_balanced_spinor}
\end{equation}
The relative phase is therefore an additional coherent coordinate above the same scalar half point.

The primitive output packet is
\begin{equation}
\boxed{
0
\prec_{\rm dep}
\mathcal P
\prec_{\rm dep}
\{N,S\}
\prec_{\rm dep}
\frac12
\prec_{\rm dep}
\ln2
\prec_{\rm dep}
\mathbb C^2.
}
\label{eq:v12_primitive_packet}
\end{equation}

\section{Branch separation}

Projectivization of the two-state carrier and the spatial geometry developed in Parts I--II proceed from the common primitive packet. The three-flavour/Standard-Model branch proceeds from the same root after the geometric carrier is established. The sibling time programme receives the primitive packet through its own typed interface.

The dependency architecture is therefore
\[
\boxed{
\mathcal C_0
\prec_{\rm dep}
\left\{
\mathcal G_X^{\rm TIR},
\mathcal B_{\rm SM}^{\rm TIR},
\mathcal B_T^{\rm IDT}
\right\},
}
\]
where $\mathcal C_0$ denotes the admitted primitive information/quantum packet.

The spatial and temporal branches later meet through the declared closure surface
\begin{equation}
\mathfrak S_X^{\rm TIR}\otimes\mathfrak T^{\rm IDT}
\longrightarrow
\mathfrak M_{XT},
\end{equation}
and the matter branch then couples to the resulting spacetime carrier. The formal completion order of these joins is recorded in Chapter~\ref{ch:v12_completion_frontier}.

\section{First-chapter invariant}

The foundational narrative used by the rest of v12 is therefore
\begin{equation}
\boxed{
\text{zero distinction}
\to\text{minimal carrier}
\to\text{nontrivial exchange}
\to\text{two poles}
\to\frac12
\to\ln2
\to\mathbb C^2.
}
\end{equation}
The next chapter resolves the geometry and information content of the half coordinate before any flavour normalization is introduced.
